\documentclass[11pt,a4paper]{article}
\usepackage[utf8]{inputenc}
\usepackage[T2A]{fontenc}
\usepackage[ukrainian,english]{babel}
\usepackage{amsmath}
\usepackage{amssymb}
\usepackage{listings}
\usepackage{graphicx}
\usepackage{hyperref}
\usepackage{geometry}
\geometry{margin=2.5cm}

\title{Адаптація Zen Crypted Dharma DSL до \\ нотації Message Sequence Chart (ITU-T/IETF Z.120)}
\author{На основі аналізу репозиторію synrc/chat}
\date{\today}

\begin{document}

\maketitle

\begin{abstract}
У цьому документі проводиться аналіз мови опису архітектури ядра чат-системи з репозиторію \texttt{synrc/chat} (файл \texttt{ARCH-KERNEL-MAPPINGS.md}), яка базується на алгебраїчних типах, морфізмах та чіткому розділенні \texttt{Action}, \texttt{Event}, \texttt{Observation} та \texttt{Predicate}. 

Ми адаптуємо цю модель до стандартизованої нотації \textbf{Message Sequence Chart} (MSC) за рекомендацією ITU-T Z.120. Пропонується відповідність між елементами kernel (Principal, Session, Post, Mutate, Seen(EventObs(...)) тощо) та конструкціями MSC (instances, messages, actions, inline expressions alt/par/opt, conditions, HMSC). 

Така адаптація дозволяє візуально та формально описувати сценарії взаємодії в акторній моделі чату, роблячи її сумісною з традиційними методами формального опису телекомунікаційних систем. 

Ключові слова: Message Sequence Chart, Z.120, акторна модель, kernel mappings, чат-система, формальні методи.
\end{abstract}

\section{Вступ}

Сучасні розподілені чат-системи часто будуються на акторній моделі з чітким розділенням між командами (\texttt{Action}), подіями (\texttt{Event}), спостереженнями та предикатами. Репозиторій \texttt{synrc/chat} (проект N2O/SYNRC) містить мінімалістичний \textit{kernel}, де всі поверхневі DSL-команди зводяться до канонічних морфізмів над типами \texttt{Principal}, \texttt{Session}, \texttt{Feed}, \texttt{Message} тощо.

Документ \texttt{ARCH-KERNEL-MAPPINGS.md} явно розмежовує:
\begin{itemize}
    \item \textbf{Action} — те, що виконується (Post, Mutate, MarkRead, Replay, View);
    \item \textbf{Event} — те, що відбулося;
    \item \textbf{Observation} — те, що спостерігається (наприклад, EventObs);
    \item \textbf{Predicate} — те, що перевіряється (Seen(...), expect).
\end{itemize}

\textbf{MSC (Z.120)} є ідеальним кандидатом для візуалізації таких мапінгів, оскільки:
\begin{itemize}
    \item Фокусується на послідовностях обміну повідомленнями між instances;
    \item Підтримує \texttt{alt}, \texttt{par}, \texttt{opt}, \texttt{loop};
    \item Дозволяє виражати actions як \texttt{out}/\texttt{in} messages або локальні \texttt{action};
    \item \texttt{expect ...} природно відображається через conditions або predicates у розширеннях MSC.
\end{itemize}

Метою роботи є створення формальної відповідності та демонстрація адаптованих сценаріїв у текстовій нотації MSC/PR.

\section{Аналіз мови kernel з synrc/chat}

\subsection{Основні об'єкти (Types)}
\begin{itemize}
    \item \texttt{Principal}, \texttt{Session}, \texttt{Feed} (Private, Group), \texttt{Message}
    \item \texttt{Event} (Typing, Presence тощо)
    \item \texttt{Observation} = EventObs(...)
    \item \texttt{Predicate} = Seen(Observation)
    \item \texttt{Action} = Post | Mutate | MarkRead | Replay | View | ...
\end{itemize}

\subsection{Ключові морфізми (Actions)}
Канонічні дії:
\begin{itemize}
    \item \texttt{Post \{session, actor, feed, payload\}}
    \item \texttt{Mutate \{session, actor, target, op\}}
    \item \texttt{MarkRead \{session, actor, boundary\}}
    \item \texttt{Replay \{session, actor, feed, after\}}
    \item \texttt{View \{session, actor, kind, ...\}}
\end{itemize}

\texttt{expect event typing bob1} зводиться не до Action, а до \texttt{Predicate(Seen(EventObs(SessionPresence\{...\}, Typing)))}.

\section{Відповідність між Kernel та MSC (Z.120)}

\begin{table}[h]
\centering
\begin{tabular}{|l|l|}
\hline
\textbf{Елемент Kernel} & \textbf{Конструкція MSC (Z.120)} \\
\hline
Session / Actor / Principal & Instance (lifeline) \\
Post, Mutate, MarkRead, View, Replay & out Message to ... (або action) \\
Event & Message label або inline action \\
Observation (EventObs) & Message з параметрами \\
Predicate Seen(...) / expect & Condition або inline opt/alt з predicate \\
State transition & Message exchange + local action \\
Alternative behaviors & alt ... endalt \\
Parallel events & par ... endpar \\
Sequence of actions & seq (implicit top-to-bottom) \\
Replay after snapshot & loop або MSC reference з after-condition \\
\hline
\end{tabular}
\caption{Відповідність елементів kernel та MSC}
\end{table}

Instances у MSC пропонуються такі:
\begin{itemize}
    \item \texttt{AliceSession}, \texttt{BobSession}
    \item \texttt{ServerKernel} або \texttt{FeedManager}
    \item \texttt{Environment} (для external predicates)
\end{itemize}

\section{Адаптовані приклади в нотації MSC/PR (Z.120)}

\subsection{Приклад 1: Надсилання повідомлення + спостереження typing}

\begin{lstlisting}[language={},basicstyle=\ttfamily\small]
msc SendMessageAndTyping;
  instance AliceSession;
    out Post {
      feed = Private(alice, bob);
      payload = {body = "hi"}
    } to BobSession;
  endinstance;

  instance BobSession;
    in Post from AliceSession;
    action ProcessMessage;
    out EventObs(SessionPresence{actor=bob1, kind=Typing}) to AliceSession;
  endinstance;

  instance AliceSession;
    in EventObs from BobSession;
    condition Seen(EventObs(Typing));
  endinstance;
endmsc;
\end{lstlisting}

\subsection{Приклад 2: Edit message (Mutate)}

\begin{lstlisting}[language={},basicstyle=\ttfamily\small]
msc EditMessageScenario;
  instance AliceSession;
    out Mutate {
      target = ExistingMessage{feed=Private(alice,bob), id="msg-123"};
      op = ReplacePayload{body="v2"}
    } to ServerKernel;
  endinstance;

  instance ServerKernel;
    in Mutate from AliceSession;
    alt
      out MutationApplied to AliceSession;
    else
      out MutationFailed to AliceSession;
    endalt;
  endinstance;
endmsc;
\end{lstlisting}

\subsection{Приклад 3: Expect + Replay after snapshot}

\begin{lstlisting}[language={},basicstyle=\ttfamily\small]
msc ReplayAfterSnapshot;
  instance AliceSession;
    out Replay {
      feed = Private(alice, bob);
      after = Some(AfterFeedSnapshot(id))
    } to ServerKernel;
  endinstance;

  instance ServerKernel;
    in Replay from AliceSession;
    loop <0, inf>
      out Event to AliceSession;
    endloop;
  endinstance;

  instance AliceSession;
    condition Satisfies(Predicate(Seen(EventsAfterSnapshot)));
  endinstance;
endmsc;
\end{lstlisting}

\section{Переваги адаптації до Z.120}

\begin{itemize}
    \item \textbf{Візуальна ясність} — сценарії стають зрозумілими для інженерів без знання OCaml.
    \item \textbf{Формальна семантика} — можна використовувати Annex B Z.120 для верифікації.
    \item \textbf{Інтеграція з TTCN-3} — сценарії MSC легко перетворюються на тести.
    \item \textbf{Масштабування} — High-level MSC (HMSC) дозволяє описувати повний життєвий цикл чату.
    \item Сумісність з SDL (Z.100) для повної специфікації стану.
\end{itemize}

\section{Висновки та подальша робота}

Адаптація kernel-мовы з \texttt{synrc/chat} до MSC Z.120 показала високу природність відповідності. Акторні взаємодії, розділення Action/Event/Observation та предикати \texttt{expect} добре відображаються через instances, messages, actions та inline expressions.

Подальша робота може включати:
\begin{itemize}
    \item Автоматичну генерацію MSC з kernel-описів.
    \item Формальну семантику мапінгу (розширення Annex B).
    \item Інтеграцію з інструментами (msc2svg, Telelogic Tau, тощо).
    \item Порівняння з UML Sequence Diagrams.
\end{itemize}

\bibliographystyle{plain}
\begin{thebibliography}{9}
\bibitem{itu-z120} ITU-T Recommendation Z.120 (02/2011), Message Sequence Chart (MSC).
\bibitem{synrc} synrc/chat repository, ARCH-KERNEL-MAPPINGS.md, 2025--2026.
\bibitem{annexb} ITU-T Z.120 Annex B (04/1998), Formal semantics of Message Sequence Charts.
\end{thebibliography}

\end{document}